\chapter{Testes e Resultados} \label{chap:testsResults}
	\section{Procedimento 01}
	\label{chap:testsResults:sec:Experimento01}

	\section{Fase 00 --- Ranking por grau de verdade paraconsistente}
	\label{chap:testsResults:sec:Fase00}

	As tabelas a seguir apresentam todas as configurações avaliadas na fase 00
	do Experimento 05, ordenadas de forma decrescente pelo grau de verdade
	paraconsistente médio ($\bar{d}$, Seção~\ref{chap:propApproach}). Aqui, $d$
	denota o grau de verdade paraconsistente de uma única repetição --- a mesma
	grandeza chamada $D_{\text{verdade}}$ no Capítulo~\ref{chap:propApproach} e
	$D_{1,0}$ na Seção~\ref{sec:conceitos} (a distância ao vértice \textit{Verdade}
	do plano paraconsistente) --- calculado a partir dos três valores de
	$\alpha$ (grau de crença) e $\beta$ (grau de descrença) obtidos nas
	repetições de cada configuração; $\bar\alpha$ e $\bar\beta$ são as médias de
	$\alpha$ e $\beta$ nessas mesmas repetições, e $\sigma_d$ é o desvio padrão
	populacional de $d$ entre as três repetições, indicando a estabilidade do
	grau de verdade da configuração. As
	Tabelas~\ref{tab:phase00ae_eeg} e~\ref{tab:phase00ae_voice} comparam as
	variantes de \textit{autoencoder} (ANN-AE e SNN-AE nas codificações direta,
	de latência e de Poisson) para EEG e voz, respectivamente. As
	Tabelas~\ref{tab:phase00hc_eeg} e~\ref{tab:phase00hc_voice} comparam as
	configurações de características manuais (família \textit{wavelet}
	$\times$ característica espectral $\times$ categoria) para EEG e voz.
	Os dados das tabelas são gerados automaticamente pelo script
	\texttt{software/nn/scripts/pipeline/\allowbreak e05\_build\_phase00\_paraconsistent\_tables.py}
	a partir dos arquivos \texttt{results/phase00/*\_summary.json}, e devem ser
	regenerados sempre que a fase 00 for reexecutada; o script escreve apenas
	os arquivos \texttt{.csv} em \texttt{tables/}.

	% Cria um longtable a partir de um CSV com colunas:
	% rank,configuration,mean_d_truth,std_d_truth,mean_alpha,mean_beta
	% Cabeçalho e zebra striping vêm de chapters/data/visualStyle.tex (SSOT).
	\newcommand{\Phasezerotable}[3]{%
		\tableZebraStripe
		\begin{longtable}{@{}rlrrrr@{}}
			\caption{#2} \label{#3} \\
			\tableHeaderRow \textbf{\#} & \textbf{Configuração} & \textbf{$\bar{d}$} & \textbf{$\sigma_d$} & \textbf{$\bar{\alpha}$} & \textbf{$\bar{\beta}$} \\
			\endfirsthead
			\tableHeaderRow \textbf{\#} & \textbf{Configuração} & \textbf{$\bar{d}$} & \textbf{$\sigma_d$} & \textbf{$\bar{\alpha}$} & \textbf{$\bar{\beta}$} \\
			\endhead
			\endfoot
			\bottomrule
			\endlastfoot
			\csvreader{#1}{}%
				{\csvcoli & \csvcolii & \csvcoliii & \csvcoliv & \csvcolv & \csvcolvi \\}%
		\end{longtable}%
	}

	\Phasezerotable{tables/phase00_ae_eeg.csv}%
		{Configurações de autoencoder (EEG) ordenadas por grau de verdade $d$}%
		{tab:phase00ae_eeg}
	\Phasezerotable{tables/phase00_ae_voice.csv}%
		{Configurações de autoencoder (voz) ordenadas por grau de verdade $d$}%
		{tab:phase00ae_voice}
	\Phasezerotable{tables/phase00_hc_eeg.csv}%
		{Configurações de características manuais (EEG) ordenadas por grau de verdade $d$}%
		{tab:phase00hc_eeg}
	\Phasezerotable{tables/phase00_hc_voice.csv}%
		{Configurações de características manuais (voz) ordenadas por grau de verdade $d$}%
		{tab:phase00hc_voice}

	\section{Fase 01 --- Autenticação DSNN}
	\label{chap:testsResults:sec:Fase01}

	A Fase 01 treina a \ac{DSNN} de autenticação sobre a extração de
	características vencedora da Fase 00 (Seção~\ref{chap:testsResults:sec:Fase00}:
	\textit{handcrafted}, \textit{wavelet} Haar, LFCC, categoria c1, para EEG e
	voz) e mede \ac{EER} e \ac{AUC} por dobra de validação cruzada. A Tabela~\ref{tab:phase01auth}
	cobre as 32 configurações do experimento --- produto cartesiano de fonte
	(EEG, voz, fusão precoce, fusão tardia) $\times$ modo de texto (dependente,
	independente) $\times$ esquema de validação cruzada ($k$-fold plano,
	$k$-fold aninhado) $\times$ padronização de características (bruto,
	padronizado) --- ordenadas de forma crescente por \ac{EER} médio (menor é
	melhor) sobre as três repetições de cada configuração. Os dados são gerados
	automaticamente pelo script
	\texttt{software/nn/scripts/pipeline/\allowbreak e05\_build\_phase01\_auth\_tables.py}
	a partir dos arquivos \texttt{results/thesis/phase01/*\_summary.json}, e
	devem ser regenerados sempre que a fase 01 for reexecutada.

	A melhor configuração (fusão precoce, dependente de texto, $k$-fold plano,
	características brutas) atinge $\text{EER} = 0{,}4534$ e
	$\text{AUC} = 0{,}5452$ --- próximo do nível de acaso
	($\text{EER} = 0{,}5$, $\text{AUC} = 0{,}5$ correspondem a um classificador
	aleatório) e, portanto, ainda distante do desempenho necessário para uma
	autenticação prática. Três padrões emergem da tabela completa: (i) EEG e as
	fusões superam voz isoladamente nas configurações mais bem ranqueadas; (ii)
	a validação cruzada em $k$-fold plano supera a aninhada dentro de cada
	combinação de fonte/modo de texto; e (iii) a padronização de características
	piora sistematicamente o \ac{EER} nesta fase, apesar de elevar o \ac{AUC} ---
	uma divergência entre as duas métricas que sugere um limiar de decisão
	inadequado sob padronização, e não necessariamente uma perda real de poder
	discriminativo, ficando como investigação futura.

	\newcommand{\Phaseonetable}[3]{%
		\tableZebraStripe
		\begin{longtable}{@{}rlrrrr@{}}
			\caption{#2} \label{#3} \\
			\tableHeaderRow \textbf{\#} & \textbf{Configuração} & \textbf{EER} & \textbf{$\sigma_{\text{EER}}$} & \textbf{AUC} & \textbf{$\sigma_{\text{AUC}}$} \\
			\endfirsthead
			\tableHeaderRow \textbf{\#} & \textbf{Configuração} & \textbf{EER} & \textbf{$\sigma_{\text{EER}}$} & \textbf{AUC} & \textbf{$\sigma_{\text{AUC}}$} \\
			\endhead
			\endfoot
			\bottomrule
			\endlastfoot
			\csvreader{#1}{}%
				{\csvcoli & \csvcolii & \csvcoliii & \csvcoliv & \csvcolv & \csvcolvi \\}%
		\end{longtable}%
	}

	\Phaseonetable{tables/phase01_auth.csv}%
		{Configurações da Fase 01 (autenticação DSNN) ordenadas por EER médio}%
		{tab:phase01auth}